\chapter{Ethics and Data Handling}
\label{app:ethics}

\section{Context and partnership}

This project was carried out in partnership with Tennor, which operates a French medical
voice agent that patients call by telephone. The work builds a real-time end-of-utterance
detector for that agent. Tennor is the data controller and operates a GDPR-compliant,
ISO/IEC 27001-certified environment. I processed the data inside Tennor's infrastructure as
part of the partnership. No participants were recruited and there was no interaction with
patients: the work reuses signals already generated by the live system.

\medskip
\noindent\fbox{\parbox{0.95\textwidth}{\small
\textbf{Correction to the statement submitted with the corpus.} The ethics statement
submitted alongside the code corpus on 21 August 2026 describes the capture as limited to
\emph{internal test organisations}. That is narrower than what actually happened: capture
was feature-flag gated at organisation level on \emph{production} organisations. This
appendix supersedes that statement and gives the accurate position. The correction was
raised with my supervisor before writing. The corpus itself can no longer be amended.}}

\section{Data used}

\textbf{Public benchmark, no personal data.} Every headline result reported in
Chapter~\ref{ch:results} runs on the LiveKit eot-bench dataset (CC-BY-4.0), a public French
turn-taking benchmark containing no Tennor data and no personal data. The benchmark
deliberately carries the weight of the evaluation, so that the method can be demonstrated
without personal data at all (Section~\ref{sec:datastrategy}).

\textbf{Production in-domain set.} A secondary, in-domain set was derived from calls to the
Tennor agent: 709 calls placed on a single production organisation, of which roughly four in five
are genuine inbound calls and the remainder calls made by internal testers to the same line.
It is used only as an in-domain validation set, and it is not included in the corpus.

\section{Legal basis and agreement with the controller}

Tennor is the data controller. Every use of Tennor data in this project was agreed with the
company before it took place, and the company reviewed the scope of the capture, the fields
retained, this appendix, and the dissertation itself. Tennor is aware of what the work covers
and of what it implies for the production agent. Nothing was collected, processed, or released
outside that agreement, and the partnership was conducted transparently on both sides.

One consequence follows directly. The in-domain production set described above, referred to as
the \emph{gold} set in the submitted corpus, cannot be shared or made public. The reason is data
sovereignty: these are French medical call recordings and transcripts held in Tennor's own
infrastructure, and they stay there. The set is therefore described and measured in this
dissertation but not distributed with it, which is why the public benchmark carries every
headline result (Section~\ref{sec:datastrategy}).

\section{Pseudonymisation, not anonymisation}

Names, first names and ages are masked in the retained transcripts, and the export is
reduced to the fields the task uses. The data nonetheless remains personal data under the GDPR: the caller-side audio is
retained and every record is linked to a production call identifier, so re-identification by
the controller remains possible. It is treated as personal data throughout, and describing
it as anonymised would be incorrect. The fields the task uses are the caller text, the recogniser's interim trajectory, the pause
and decision timestamps, an audio reference, and the resume and interruption flags. The
export is reduced to these.

\section{Reduced human exposure}

Labels are produced automatically by weak supervision, from whether the caller resumed
speaking after a pause (Section~\ref{sec:instrumentation}). Building the dataset therefore
required nobody to sit and listen to call content, which limits human exposure to what
patients said. The cost of that choice is label noise, and it bounds the result reported in
Section~\ref{sec:gold}. It is an ethical decision with a measurable methodological price,
which is worth stating in both directions.

\section{Storage, access and the corpus}

The data and any derived audio or transcripts stay within Tennor's infrastructure. Access
was restricted to me. Nothing was copied into the repository, the corpus, or any public
location: \texttt{.gitignore} excludes call audio, transcripts and model weights, and the
corpus was packaged and checked so that none of those files were included. The corpus
contains the source code, the dataset-generation code, the evaluation on the public dataset,
the design documents and this statement.

\section{Risks and mitigations}

\begin{table}[H]
\centering
\small
\begin{tabular}{@{}p{5.5cm}p{8.0cm}@{}}
\toprule
Risk & Mitigation \\
\midrule
Exposure of patient data & Nothing leaves Tennor's infrastructure; no audio, transcript or
derived record in the corpus; access limited to the author \\
Re-identification & Names, first names and ages masked in the retained transcripts; treated as personal data throughout rather than described as anonymised; nothing exported beyond Tennor's infrastructure \\
Accidental inclusion in the corpus & \texttt{.gitignore} rules plus packaging that excludes
audio, transcripts and weights, verified before submission \\
Over-collection & Only caller-side, task-relevant fields captured per micro-pause \\
Human exposure to call content & Automatic weak-supervision labelling instead of manual
listening \\
Clinical impact of a wrong cut & Cost-sensitive decision with the interruption cost stated
explicitly, and cut-in rates reported openly; the system augments rather than replaces
existing safeguards \\
\bottomrule
\end{tabular}
\end{table}
